\section{Билет 55. Распределение Максвелла молекул по величине скорости. Формула Максвелла в приведенном виде.}

\begin{definition}
Распределение Максвелла по модулю скорости задаётся формулой:
\begin{equation}
    \frac{dN}{N} = F(v) \, dv = 4\pi \left( \frac{m}{2\pi k T} \right)^{3/2} v^2 \exp\left( -\frac{m v^2}{2 k T} \right) dv, \label{eq:maxwell_v_55}
\end{equation}
где $dN/N$ --- относительное число молекул, модули скоростей которых лежат в интервале от $v$ до $v+dv$. Функция $F(v)$ нормирована на единицу: $\int_0^\infty F(v) \, dv = 1$.
\end{definition}

\begin{definition}[Приведённая форма распределения]
Для решения многих задач удобно использовать безразмерные (приведённые) единицы скорости. Введём относительную скорость:
\begin{equation}
    u = \frac{v}{v_{\text{вер}}}, \quad \text{где } v_{\text{вер}} = \sqrt{\frac{2kT}{m}} \text{ --- наиболее вероятная скорость.}
\end{equation}
Тогда $v = u v_{\text{вер}}$ и $dv = v_{\text{вер}} \, du$. Подставляя в \eqref{eq:maxwell_v_55} и учитывая, что $v_{\text{вер}}^2 = 2kT/m$, получаем:
\begin{equation}
    \frac{dN}{N} = \frac{4}{\sqrt{\pi}} u^2 e^{-u^2} \, du = f(u) \, du. \label{eq:reduced_maxwell}
\end{equation}
\end{definition}

\begin{property}[Универсальность приведённой формы]
Функция $f(u) = \frac{4}{\sqrt{\pi}} u^2 e^{-u^2}$ не зависит ни от рода газа ($m$), ни от температуры ($T$). Это означает, что кривые распределения Максвелла для любых газов и температур подобны друг другу и различаются только масштабом по осям.
Относительная доля молекул со скоростями в интервале от $u_1$ до $u_2$ вычисляется по формуле:
\begin{equation}
    \frac{\Delta N}{N} = \int_{u_1}^{u_2} \frac{4}{\sqrt{\pi}} u^2 e^{-u^2} \, du.
\end{equation}
Для узких интервалов $\Delta u \ll 1$ можно использовать приближение $\Delta N/N \approx f(u) \Delta u$.
\end{property}

\begin{remark}
\textbf{Характерные значения распределения:}
\begin{itemize}
    \item $u \in [0.5; 1.5]$ --- около $70.7\%$ молекул (скорости отличаются от наиболее вероятной не более чем на $50\%$);
    \item $u \in [1.5; 2.0]$ --- около $16.6\%$ молекул;
    \item $u > 3$ --- лишь $0.04\%$ молекул обладают скоростью, превышающей $v_{\text{вер}}$ более чем в три раза;
    \item $u > 5$ --- вероятность ничтожно мала ($\sim 8 \cdot 10^{-9}$).
\end{itemize}
\textbf{Пример:} Доля молекул со скоростями, отличающимися от $v_{\text{вер}}$ не более чем на $2\%$ ($u \in [0.98; 1.02]$, $\Delta u = 0.04$), равна $\Delta N/N \approx f(1) \cdot 0.04 = \frac{4}{\sqrt{\pi}} e^{-1} \cdot 0.04 \approx 0.033$, или $3.3\%$.
\end{remark}
